% Independent full coefficient-divisor derivation. Prefix HFD.
% Requires amsmath, amssymb, hyperref.
\section{The heat divisor before coefficient-support synchronization}
\label{sec:coefficient-flag-heat-divisor}

\paragraph{Sources and original objects.}
The full coefficient algebra is SHP1--28 of
\nolinkurl{SUPPORTED_HEAT_PRODUCT.tex}, with analytic completion
HSL15--35 of \nolinkurl{FULL_LATTICE_HEAT_OPERATOR.tex}.
The original foundation is retained under the byline The Clankers in
\href{https://github.com/KokunoYumeto/zeta-function-research-reader/blob/a2851f9eb352e7e4376bc629de86fa4ed9c1c434/workbenches/splitzero-tandem/foundations/split-support-geometry-v11/split_support_geometry_arithmetic_curve_v11.tex}{\emph{Split
Support Geometry: Universal Zero Fibres, Arithmetic Curves, Frobenius
and Perfect Quantization}, v11, Definition \texttt{def:lattice-split}}.
The actual analytic kernel is that of Brad Rodgers and Terence Tao,
\href{https://arxiv.org/abs/1801.05914v5}{\emph{The De Bruijn--Newman
constant is non-negative}, arXiv:1801.05914v5}.
The complete proofs THS8--15 and THS18--23 in
\nolinkurl{THETA_HEAT_FULL_SUPPORT.tex} establish its coefficient
formula, nonvanishing of every even coefficient at real time,
membership in the stated analytic space and the exact heat group.
No theorem about its zero locations is used here.

The prime receiver read for this calculation is HFP11--12 and
HFP34--37 of \nolinkurl{HEAT_FLAG_PRIME_SPECTRUM.tex}, lines
216--273 and 682--787, at source SHA--256
\begin{center}\small
\nolinkurl{e51115eaa2e70bb97cdf32d113222eace1e37260e1ec6781dbefe62517dab63f}.
\end{center}
The synchronized divisor comparison is HZD1--8 of
\nolinkurl{FULL_LATTICE_HEAT_DIVISOR.tex}, read at source
\begin{center}\small
\nolinkurl{986b34093592f4949645a02a2413817809476cb82519f46f61471ba1a03ede3a}.
\end{center}
The algebra below calculates the congruence on the full
coefficient flags before that synchronization. All receiver
maps used from these sources are proved again in the required
quotient coordinates.

\subsection{The coefficient fibres and the supported equation}

Let $L$ be the entire nontrivial bounded distributive lattice,
and retain the original Taylor coefficient space
\[
 \mathcal E=\left\{f(z)=\sum_{n\ge0}a_nz^n:
       \sum_{n\ge0}|a_n|\sqrt{n!}\,h^n<\infty
                         \text{ for every }h>0\right\}.
\]
Equivalently, $f$ is entire and
$p_\alpha(f)=\sup_z|f(z)|e^{-\alpha|z|^2}<\infty$ for
every $\alpha>0$, by HSL16. It is a ring under ordinary
function multiplication, with the proved bound
$\|fg\|_h\le\|f\|_{\sqrt2h}\|g\|_{\sqrt2h}$.
Write $\mathcal E^{(0)}$ and $\mathcal E^{(1)}$ for its
even and odd subspaces. Their projections select the even
and odd Taylor coefficients, respectively, and preserve
every defining seminorm.

Let $\mathcal F(L)$ be the full flag semiring. Its elements
are sequences $\nu_i\in L$ satisfying $\nu_i\ge\nu_{i+2}$.
Addition is coordinate join and multiplication is
\begin{equation}\tag{HFD1}
 (\nu\circledast\mu)_0
    =(\nu_0\wedge\mu_0)\vee(\nu_1\wedge\mu_1),\qquad
 (\nu\circledast\mu)_n
    =\bigvee_{r=0}^{n}\nu_r\wedge\mu_{n-r}\quad(n\ge1).
\end{equation}
The unit is $\mathbf u=(1_L,0_L,0_L,\ldots)$.
SHP3--9 prove its algebra laws by the exact
parity-tail closure identity. The coefficient algebra is
\begin{equation}\tag{HFD2}
 \mathcal A_L=\{(f,\nu):\nu\in\mathcal F(L),\
       f=\sum a_nz^n\in\mathcal E,\
       a_n\ne0\Longrightarrow\nu_n=1_L\},\qquad
 (f,\nu)\circ(g,\mu)=(fg,\nu\circledast\mu).
\end{equation}
The zero is $(0,\mathbf0)$ and the unit is $(1,\mathbf u)$.
Every label in $L$ remains available at a zero coefficient.

For each flag retain its complete amplitude fibre and its
two full-parity indicators:
\begin{equation}\tag{HFD3}
 \begin{gathered}
 \mathcal E_\nu
    =\{f=\sum a_nz^n\in\mathcal E:
                       \nu_n\ne1_L\Longrightarrow a_n=0\},\\
 \epsilon_p(\nu)=
  \begin{cases}1,&\nu_{2j+p}=1_L\text{ for every }j\ge0,\\
                0,&\text{otherwise},
  \end{cases}
       \qquad(p=0,1),\\
 V_\nu=\epsilon_0(\nu)\mathcal E^{(0)}
                  \oplus\epsilon_1(\nu)\mathcal E^{(1)}.
 \end{gathered}
\end{equation}
These are complex vector spaces. For a parity with
$\epsilon_p(\nu)=0$, its component of $\mathcal E_\nu$
consists of polynomials: choose an index of that parity
whose label is not top. Every later label of that parity
is below it, and hence also not top, so all those
coefficients vanish. Before that index only the initial
top-labelled coefficients are allowed. For a parity
with indicator $1$, the entire corresponding parity
subspace of $\mathcal E$ is allowed. This is an exact
description of $\mathcal E_\nu$, not a change to the
remaining non-top labels in $\nu$.

Let $s\in\mathbb C$ be the actual theta time, and set
\begin{equation}\tag{HFD4}
 H=H_s,\qquad
 \eta_{2j}=1_L,\quad\eta_{2j+1}=0_L,\qquad
 h_s=(H_s,\eta),\quad z_\eta=(0,\eta).
\end{equation}
The analytic function $H_s$ is even, belongs to $\mathcal E$,
and is not zero, by the actual heat group and the positive
real kernel at time zero.
Both elements in (HFD4) belong to the original coefficient
algebra. The supported equation is $h_s\sim z_\eta$,
not $h_s\sim(0,\mathbf0)$.

The crucial support multiplication is exactly
\begin{equation}\tag{HFD5}
 \eta\circledast\nu=\kappa(\nu),\qquad
 \kappa(\nu)_{2j}=\nu_0,\quad
 \kappa(\nu)_{2j+1}=\nu_1 .
\end{equation}
At an even positive output index in (HFD1), the term
with the $\eta$ index equal to that output index
contributes $\nu_0$, and the other contributing labels
are at most $\nu_0$. At an odd index the largest
contributing label is $\nu_1$. At degree zero (HFD1)
gives exactly $\nu_0$. This proves every coordinate.

\subsection{The exact generated congruence}

The congruence generated by the single supported equation
in (HFD4) is the following explicit equivalence:
\begin{equation}\tag{HFD6}
 (f,\nu)\equiv_s(g,\mu)
 \quad\Longleftrightarrow\quad
 \nu=\mu\ \text{ and }\ f-g=H_s b
                        \text{ for some }b\in V_\nu .
\end{equation}
In particular every original coefficient label is
unchanged in an equivalence class.

Here is the full congruence proof.
First $H_sV_\nu\subset\mathcal E_\nu$, since $H_s$
is even and multiplication preserves parity and
$\mathcal E$. On a permitted parity the flag is
top in every degree; no other parity is produced.
Thus (HFD6) is an equivalence on each vector-space
fibre. Under addition of a context with flag $\mu$,
$V_\nu\subset V_{\nu\vee\mu}$, so the equivalence is
preserved by addition.

For multiplication the needed exact fact is
\begin{equation}\tag{HFD7}
 b\in V_\nu,\quad k\in\mathcal E_\mu
       \quad\Longrightarrow\quad
 bk\in V_{\nu\circledast\mu}.
\end{equation}
Decompose both functions into their even and odd
parts. If the parity-$p$ part of $b$ is nonzero,
then every parity-$p$ label of $\nu$ is top.
If the parity-$q$ part of $k$ is nonzero, some
Taylor coefficient in that parity is nonzero, so
its label and therefore $\mu_q$ are top.
The output flag has every parity-$(p+q\bmod2)$
label top. For an odd output index this follows
by using the degree-$0$ or degree-$1$ label of
$\mu$ in (HFD1) and the appropriate top label
of $\nu$. The same argument works for positive
even indices. For output degree zero, the
even--even case uses $\nu_0\wedge\mu_0$,
and the odd--odd case uses the indispensable
additional term $\nu_1\wedge\mu_1$.
Thus each nonzero parity product in $bk$ lies
in an allowed full parity of the output.
Joining these possibilities proves (HFD7).
It follows that multiplying the difference $H_sb$
by any context amplitude again has the form in
(HFD6) for the product flag. This proves
multiplicative compatibility. Compatibility for
two equivalent factors follows by changing one
factor at a time. Hence (HFD6) is a semiring
congruence, and it contains the generating pair
because $1\in V_\eta$.

Conversely every pair in (HFD6) is produced by
one elementary use of that generating pair.
For $b=b_0+b_1\in V_\nu$, let $\chi_b$ be the
parity-constant flag that is top on parity $p$
when $b_p\ne0$, and bottom otherwise. Then
$(b,\chi_b)\in\mathcal A_L$,
$\chi_b\le\nu$, and
$\eta\circledast\chi_b=\chi_b$.
Multiply $h_s\sim z_\eta$ by this element and
add the context $(g,\nu)$. The two results are
\[
 (g,\nu)+(H_sb,\chi_b)=(f,\nu),\qquad
 (g,\nu)+(0,\chi_b)=(g,\nu).
\]
This proves that every equivalence in (HFD6)
is generated by the required relation. The
congruence calculation is therefore exact.

The divisor before synchronization is the concrete
semiring
\begin{equation}\tag{HFD8}
 \mathcal D^{\mathrm{flag}}_{s,L}
   =\mathcal A_L/\!\equiv_s
   =\coprod_{\nu\in\mathcal F(L)}
                \bigl(\mathcal E_\nu/H_sV_\nu\bigr)\times\{\nu\}.
\end{equation}
For representatives the operations are
\begin{equation}\tag{HFD9}
 ([f]_\nu,\nu)+([g]_\mu,\mu)
       =([f+g]_{\nu\vee\mu},\nu\vee\mu),\qquad
 ([f]_\nu,\nu)\circ([g]_\mu,\mu)
       =([fg]_{\nu\circledast\mu},\nu\circledast\mu).
\end{equation}
They are well defined by the congruence proof.
The unit and zero retain the flags $\mathbf u$
and $\mathbf0$, respectively. Formula (HFD8)
records every support-dependent fibre.
On the actual $\eta$ fibre only even multipliers
are allowed: its amplitude space is
$\mathcal E^{(0)}/H_s\mathcal E^{(0)}$.

\subsection{The actual heat function removes the comparison defect}

The following analytic fact is proved directly on
the original completion: every even, nonpolynomial
function in $\mathcal E$ has infinitely many distinct
complex zeros.
Suppose instead that such an $F$ had only finitely
many zeros, and form their monic zero polynomial $p$,
with all multiplicities. Evenness gives equal
multiplicities at $a$ and $-a$, and an even
multiplicity at zero. Thus $p$ is even.
The quotient $G=F/p$ is entire, even and zero-free.
It is still in $\mathcal E$: outside a sufficiently
large disc $|p(z)|\ge1$, and on the remaining
compact disc $G$ is bounded. Consequently every
positive Gaussian bound for $F$ gives one for $G$.

The entire function $G'/G$ has an entire primitive,
obtained by integrating its Taylor series.
Choosing its additive constant to be a logarithm
of $G(0)$ gives an entire $\ell$ with $G=e^\ell$.
For every $\alpha>0$,
\[
 \operatorname{Re}\ell(z)
   \le \log p_\alpha(G)+\alpha|z|^2 .
 \tag{HFD10}
\]
Write $\ell(z)=\sum_{n\ge0}c_nz^n$ and, on the circle
of radius $R$, put
$u(\theta)=\operatorname{Re}\ell(Re^{i\theta})
                         -\operatorname{Re}c_0$.
Its average is zero by the convergent Taylor
series, and
$u\le M_{\alpha,R}:=
\log p_\alpha(G)+\alpha R^2-\operatorname{Re}c_0$.
This upper bound is nonnegative since
$p_\alpha(G)\ge|G(0)|$.
The zero average gives
$\int_0^{2\pi}|u|=2\int_0^{2\pi}u_+
                         \le4\pi M_{\alpha,R}$.
Fourier coefficient extraction from the same
Taylor series therefore gives, for $n\ge1$,
\[
 |c_n|R^n
 =\left|\frac1\pi\int_0^{2\pi}
                  u(\theta)e^{-in\theta}\,d\theta\right|
 \le4M_{\alpha,R}.
 \tag{HFD11}
\]
For $n\ge3$, let $R\to\infty$ with $\alpha$
fixed; then $c_n=0$. For $n=2$ this yields
$|c_2|\le4\alpha$, and $\alpha\downarrow0$
gives $c_2=0$. Thus $\ell=c_0+c_1z$.
Since $G$ is even, $G'(0)=0$, forcing $c_1=0$.
It follows that $F$ is a constant times $p$,
contradicting its nonpolynomiality.
This proves the analytic fact without assuming
where any of the zeros lie.

Apply it to every actual $H_s$, including complex
time. The function $H_0$ has nonzero Taylor
coefficients in every even degree by THS9,
so it is not a polynomial.
If $H_s$ were polynomial, the identity
$H_0=T_{-s}H_s$ would make $H_0$ polynomial:
the heat series terminates on a polynomial.
Also $H_s$ is even and is in $\mathcal E$.
We have proved
\begin{equation}\tag{HFD12}
 H_s\text{ has infinitely many distinct zeros},\qquad
 H_sb\text{ polynomial},\ b\in\mathcal E
      \quad\Longrightarrow\quad b=0 .
\end{equation}
For the second assertion a polynomial $H_sb$
vanishes at every zero of $H_s$, and hence is zero.
The ring of entire functions is an integral
domain: where $H_s$ is nonzero, $b$ vanishes
on an open set, and the identity theorem gives
$b=0$ everywhere. This also proves that $(H_s)$
is a proper ideal in $\mathcal E$.

The exact comparison defect now vanishes for this
actual function:
\begin{equation}\tag{HFD13}
 H_s\mathcal E\cap\mathcal E_\nu=H_sV_\nu,\qquad
 \frac{H_s\mathcal E\cap\mathcal E_\nu}{H_sV_\nu}=0
                     \quad(\nu\in\mathcal F(L)).
\end{equation}
Only one inclusion remains to prove.
Suppose $H_sb\in\mathcal E_\nu$.
For a parity with $\epsilon_p(\nu)=0$,
that parity component is a polynomial by
(HFD3). Since $H_s$ is even, that component
is $H_sb_p$, which forces $b_p=0$ by
(HFD12). All other parity components of $b$
are allowed in $V_\nu$.
This proves the reverse inclusion and the
entire displayed defect calculation.

Thus the divisor has an injective pair description,
but with a proved, restricted image:
\begin{equation}\tag{HFD14}
 \begin{gathered}
 \mathcal D^{\mathrm{flag}}_{s,L}
      \hookrightarrow(\mathcal E/(H_s))\times\mathcal F(L),
       \qquad([f]_\nu,\nu)\longmapsto([f],\nu),\\
 \operatorname{image}
   =\{(c,\nu):c\in\operatorname{image}
                   (\mathcal E_\nu\to\mathcal E/(H_s))\}.
 \end{gathered}
\end{equation}
Injectivity is exactly (HFD13).
The description of the image follows from (HFD8)
and is both necessary and sufficient.
It is not the unrestricted product: for example
$([1],\mathbf0)$ is not in it, since
$\mathcal E_{\mathbf0}=0$ and $[1]\ne0$ by
(HFD12). No support constraint has been removed.

\subsection{The induced complete coefficient quotient}

Retain $K_L=K\otimes_{\mathbb B}L$ and the
injective label map $\iota_L(\lambda)=1_K\otimes\lambda$,
with retraction $\rho_L$.
The exact full coefficient map is
$\boldsymbol\Phi(f,\nu)
   =(f,(\iota_L(\nu_i))_i)$.
In its full target
$\mathcal E\times\mathcal F(K_L)$, the congruence
generated by $\boldsymbol\Phi(h_s)\sim
\boldsymbol\Phi(z_\eta)$ is exactly
\begin{equation}\tag{HFD15}
 (f,w)\sim(g,v)
   \quad\Longleftrightarrow\quad
   w=v,\quad f-g\in H_s\mathcal E .
\end{equation}
Reduction of the first coordinate modulo $(H_s)$
is a homomorphism coequalizing the pair, so
it gives one containment of congruences.
For the other, multiply that pair in the full
target by $(b,\mathbf0)$ and then add $(g,w)$.
This gives $(g+H_sb,w)\sim(g,w)$.
The full target really permits this multiplier
even when $b\ne0$; the original coefficient
carrier does not. This is why its restriction
was calculated in (HFD13), rather than assumed.

The complete quotient square is
\begin{equation}\tag{HFD16}
 \begin{array}{ccc}
 \mathcal A_L&\xrightarrow{\ \boldsymbol\Phi\ }&
               \mathcal E\times\mathcal F(K_L)\\
 \big\downarrow&&\big\downarrow\\
 \mathcal D^{\mathrm{flag}}_{s,L}&
 \xrightarrow{\ \boldsymbol\Phi_s\ }&
               (\mathcal E/(H_s))\times\mathcal F(K_L),
 \end{array}
 \qquad
 \boldsymbol\Phi_s([f]_\nu,\nu)
         =([f],(\iota_L(\nu_i))_i).
\end{equation}
The lower map is injective by (HFD13) and the
label retraction. Its image is exactly the
support-constrained image in (HFD14) after
applying $\iota_L$ to each flag coordinate.
Both composites give the displayed pair.
The square is a pushout: compatible maps from
the upper-right and lower-left objects impose
exactly the generating relation on the upper-right
object, whose entire congruence was calculated
in (HFD15). The resulting factorization is unique.
This proves the complete-map statement including
the induced relation and its exact restriction.

Let $J=(1,(1_L,1_L,\ldots))$ be the synchronization
idempotent of SHP20. Its image $\overline J$ in
the divisor still acts by
\begin{equation}\tag{HFD17}
 \overline J\circ([f]_\nu,\nu)
       =([f]_{\mathrm{const}(\nu_0\vee\nu_1)},
                        \mathrm{const}(\nu_0\vee\nu_1)),\qquad
 \mathcal D^{\mathrm{flag}}_{s,L}[\overline J^{-1}]
                 \simeq G_L(\mathcal E/(H_s)).
\end{equation}
For a constant flag of label $\lambda\ne1_L$,
the only amplitude is zero. For the constant
top flag its amplitude quotient is
$\mathcal E/(H_s)$, since both parity indicators
in (HFD3) are one. These are exactly the
carrier and operations of the right side.
Multiplication by an idempotent realizes its
localization, with the idempotent as internal
unit, proving the isomorphism and all its
coordinates. In particular the map is
$([f]_\nu,\nu)\mapsto([f],\nu_0\vee\nu_1)$.
This is the exact connection to the earlier
synchronized supported divisor.

\subsection{The mixed prime really descends before synchronization}

The quotient does not identify different flags,
so it retains the full head-pair receiver
\begin{equation}\tag{HFD18}
 \begin{gathered}
 \overline h:\mathcal D^{\mathrm{flag}}_{s,L}
          \longrightarrow L[C_2],\qquad
 \overline h([f]_\nu,\nu)=(\nu_0,\nu_1),\\
 (a,b)+(c,d)=(a\vee c,b\vee d),\\
 (a,b)(c,d)=((a\wedge c)\vee(b\wedge d),
                     (a\wedge d)\vee(b\wedge c)),\\
 0=(0_L,0_L),\qquad1=(1_L,0_L).
 \end{gathered}
\end{equation}
The two head formulas in (HFD1) prove that
this is a unital homomorphism.
It is onto: use zero amplitude and the
periodic flag $(a,b,a,b,\ldots)$.
Moreover $z_\eta^2=z_\eta$, and its image
$\overline z_\eta=\overline h_s$ acts by
zeroing the amplitude and making the two
heads periodic, by (HFD5). Thus
\begin{equation}\tag{HFD19}
 \mathcal D^{\mathrm{flag}}_{s,L}
                  [\overline z_\eta^{-1}]
       \simeq L[C_2],\qquad
 ([f]_\nu,\nu)\longmapsto(\nu_0,\nu_1).
\end{equation}
The inverse corner map is
$(a,b)\mapsto([0],(a,b,a,b,\ldots))$.
It is unital into the corner whose identity
is $\overline z_\eta$. Distinct flags have
not been identified by the quotient, so this
map is injective as well as onto that corner.
This proves the localization directly.
Under the complete map (HFD16), the head
receiver commutes with
$L[C_2]\hookrightarrow K_L[C_2]$ given by
$\iota_L$ on both coordinates; it has the
coordinate retraction $\rho_L$ on both.

Let $\mathfrak q$ be any proper lattice prime
and $\chi_{\mathfrak q}:L\to\mathbb B$ its
character. It preserves join and meet:
membership of a join in the ideal is equivalent
to membership of both summands, and primality
tests a meet. The full receiver above gives
\begin{equation}\tag{HFD20}
 \overline\pi_{\mathfrak q}([f]_\nu,\nu)
   =\{\bar p:p=0,1,\ \nu_p\notin\mathfrak q\}
                         \in\mathscr P(C_2).
\end{equation}
The target has union as addition and group-set
addition as multiplication, zero $\varnothing$
and unit $\{\bar0\}$. Formula (HFD18) proves
that this is a surjective unital homomorphism.
The subset
$I_{\mathrm{mix}}=\{\varnothing,C_2\}$
is an ordinary prime ideal: it is closed under
union and multiplication by every subset, and
its complement consists of the two singletons,
whose products are again singletons. Therefore
\begin{equation}\tag{HFD21}
 \overline P^{\mathrm{mix}}_{\mathfrak q}
  =\overline\pi_{\mathfrak q}^{-1}(I_{\mathrm{mix}})
  =\{([f]_\nu,\nu):
       \chi_{\mathfrak q}(\nu_0)
                    =\chi_{\mathfrak q}(\nu_1)\}
\end{equation}
is a proper ordinary prime of the divisor.

This also proves the required saturation of
the source prime, rather than assuming it.
Every pair in the generated congruence (HFD6)
has the same complete flag, and hence the same
$\pi_{\mathfrak q}$ value. Membership of its
inverse image of $I_{\mathrm{mix}}$ is therefore
constant on each congruence class.
Its quotient image is exactly (HFD21).
Addition and multiplication in the quotient
preserve that ideal, and primality follows by
lifting any product to representatives and
using source primality, or directly from the
displayed receiver. The prime is not subtractive:
in the target, $C_2\in I_{\mathrm{mix}}$ and
$C_2\cup\{\bar0\}\in I_{\mathrm{mix}}$, but
$\{\bar0\}\notin I_{\mathrm{mix}}$.
The same counterexample lifts using periodic
zero-amplitude flags.

The distinction relevant to synchronization is
exact:
\begin{equation}\tag{HFD22}
 \overline J\in\overline P^{\mathrm{mix}}_{\mathfrak q},
 \qquad
 \overline h_s=\overline z_\eta
        \notin\overline P^{\mathrm{mix}}_{\mathfrak q}.
\end{equation}
Their head pairs are $(1_L,1_L)$ and
$(1_L,0_L)$, respectively.
Thus this prime survives the head localization
(HFD19), but is excluded by localization at
$\overline J$ in (HFD17).
The supported equation kills an analytic
amplitude while leaving $z_\eta$ different
from the absorbing zero; that is why the
proper prime (HFD21) can remain.

\subsection{The off-real analytic comparison retaining all flags}

Put $\Omega=\mathbb C\setminus\mathbb R$ and
$\mathcal O=\mathcal O(\Omega)$, keeping both
connected components. Restriction is a unital
ring map $\mathcal E\to\mathcal O$.
Define
\begin{equation}\tag{HFD23}
 R_s=\mathcal O/(H_s|_\Omega),\qquad
 I_s^{\mathrm{off}}
   =\{f\in\mathcal E:f|_\Omega\in
                           H_s|_\Omega\,\mathcal O\},\qquad
 B_s=\mathcal E/I_s^{\mathrm{off}}
          \hookrightarrow R_s .
\end{equation}
The last map identifies $B_s$ with the exact
image of the coefficient ring in the analytic
quotient. Its injectivity is the definition
of the kernel $I_s^{\mathrm{off}}$.
No assertion that arbitrary holomorphic
classes are represented by $\mathcal E$
is needed or made.

The ideal in (HFD23) is explicitly determined
by all off-real zero jets:
\begin{equation}\tag{HFD24}
 f\in I_s^{\mathrm{off}}
 \quad\Longleftrightarrow\quad
 f^{(j)}(\zeta)=0
 \quad\text{for every }\zeta\in\Omega
       \text{ with }H_s(\zeta)=0,\
          0\le j<\operatorname{ord}_{\zeta}H_s .
\end{equation}
Indeed near such a zero write
$H_s(z)=(z-\zeta)^m u(z)$ with $u(\zeta)\ne0$.
The quotient $f/H_s$ extends holomorphically
there exactly when the first $m$ Taylor
coefficients of $f$ vanish.
Away from the zeros it is already holomorphic.
When all the stated conditions hold the
extensions agree with that quotient on
punctured neighbourhoods and hence agree
with one another, defining an element of
$\mathcal O(\Omega)$. This proves both
directions with every multiplicity retained.

The exact off-real quotient before synchronization is
\begin{equation}\tag{HFD25}
 \begin{gathered}
 \mathcal D^{\mathrm{flag,off}}_{s,L}
  =\coprod_{\nu\in\mathcal F(L)}
       \bigl(\mathcal E_\nu/
          (I_s^{\mathrm{off}}\cap\mathcal E_\nu)\bigr)
                                               \times\{\nu\}\\
  \hookrightarrow R_s\times\mathcal F(L),\qquad
   ([f]_\nu,\nu)\longmapsto([f|_\Omega],\nu).
 \end{gathered}
\end{equation}
Its image consists exactly of pairs whose
first coordinate has a representative from
$\mathcal E_\nu$ for the displayed flag.
The equivalence on the source is equality
of flags together with the full jet test
(HFD24) on $f-g$. This is the kernel of
the actual semiring map
$\mathcal A_L\to R_s\times\mathcal F(L)$,
so its operations are (HFD9) with the
new amplitude quotients; it is a
congruence quotient, not an assumed
unrestricted product.
Since $H_sV_\nu\subset I_s^{\mathrm{off}}
\cap\mathcal E_\nu$, the source divisor
(HFD8) maps onto (HFD25).
On each fixed flag, the exact additional
identified subspace is
\[
 (I_s^{\mathrm{off}}\cap\mathcal E_\nu)/
                                     H_sV_\nu .
 \tag{HFD26}
\]
This proves the complete comparison defect
to the off-real analytic receiver.

The coefficient map retains its full form:
the map in (HFD25) followed by
$\operatorname{id}_{R_s}\times\mathcal F(\iota_L)$
is an injective map to
$R_s\times\mathcal F(K_L)$.
It agrees with the composite of (HFD16)
and the ring map
$\mathcal E/(H_s)\to R_s$.
All head maps and the prime (HFD21)
factor through this further quotient too,
because its equivalences still preserve
every label. Their properness and primality
follow by exactly the same receiver proof.

Synchronization of (HFD25) has the precise
image and connection to HZD:
\begin{equation}\tag{HFD27}
 \mathcal D^{\mathrm{flag,off}}_{s,L}
       [\overline J^{-1}]
    \simeq G_L(B_s)\hookrightarrow
                 G_L(R_s)=\mathscr D_{s,L}.
\end{equation}
The map sends a class to
$([f|_\Omega],\nu_0\vee\nu_1)$.
Its image contains every nonzero amplitude
class in $B_s$ with top label, by choosing
a representing $f$ and the constant top
flag. It contains every labelled zero by
choosing zero amplitude and a constant
flag of that label. These exhaust
$G_L(B_s)$, including when $B_s$ is the
zero ring. The idempotent corner calculation
with $J$ is unchanged, proving (HFD27).
The last inclusion is induced by the
injective ring map in (HFD23), with every
lattice label retained. This is the exact
map to HZD7 and states its possibly larger
analytic amplitude ring explicitly.

For a zero-free restriction $H_s|_\Omega$,
the reciprocal $1/H_s$ is holomorphic on
$\Omega$, so $R_s=0$ and
$I_s^{\mathrm{off}}=\mathcal E$.
Conversely, $R_s=0$ gives a holomorphic
inverse and excludes zeros by evaluation.
Accordingly the complete retained object
in that case is
\begin{equation}\tag{HFD28}
 R_s=0
 \quad\Longrightarrow\quad
 \mathcal D^{\mathrm{flag,off}}_{s,L}
          \simeq\mathcal F(L),\qquad
 G_L(R_s)\simeq L .
\end{equation}
The first equality follows in every fibre
of (HFD25), whose amplitude quotient is
then zero. It retains every flag and the
proper mixed prime (HFD21).
Thus the existence of that residual
support prime is not evidence for an
off-real zero; it persists even after
the analytic amplitude quotient vanishes.
This conclusion follows from the computed
objects and does not decide whether the
actual restriction has such a zero.

\subsection{The transported product and the correct time in the divisor}

Keep the actual theta time $s$ distinct
from the algebra transport time $t$.
On $(\mathcal A_L,\star_t)$ the full
heat map
$U_{-t}(f,\nu)=(T_{-t}f,\nu)$ is a
unital algebra isomorphism to
$(\mathcal A_L,\circ)$.
It sends $h_s$ to $h_{s-t}$ and
$z_\eta$ to itself by the proved
actual heat law. Therefore the exact
generated transported congruence is
\begin{equation}\tag{HFD29}
 (f,\nu)\equiv_{s;t}(g,\mu)
 \quad\Longleftrightarrow\quad
 \nu=\mu,\quad
 T_{-t}(f-g)=H_{s-t}b
                         \text{ for some }b\in V_\nu .
\end{equation}
This follows by applying the algebra
isomorphism to the entire proved
congruence (HFD6), not only to its
generating pair.
The spaces $\mathcal E_\nu$ and $V_\nu$
are stable under $T_t$ and $T_{-t}$:
each allowed unrestricted parity stays
in its parity, and every restricted
parity is a finite-degree polynomial
space on which the even derivatives
only lower the degree.
Thus (HFD29) also gives the exact
vector-space fibres in the original
coefficient coordinates.

Writing $\mathcal D^{\mathrm{flag},\star_t}_{s,L}$
for this quotient, the full comparison is
\begin{equation}\tag{HFD30}
 \begin{gathered}
 \mathcal D^{\mathrm{flag},\star_t}_{s,L}
      \xrightarrow[{\ [f]_\nu\mapsto[T_{-t}f]_\nu\ }]{\ \sim\ }
                   \mathcal D^{\mathrm{flag}}_{s-t,L},\\
 \mathcal D^{\mathrm{flag},\star_t}_{s,L}
       \longrightarrow G_L(\mathcal E/(H_{s-t}))
       \longrightarrow
          G_L(\mathcal O(\Omega)/(H_{s-t}|_\Omega)).
 \end{gathered}
\end{equation}
The global amplitude representative is $T_{-t}f$.
The complete off-real jet test is (HFD24)
with $H_{s-t}$ and $T_{-t}(f-g)$.
The full coefficient quotient map likewise
has amplitude ring
$(\mathcal E,\star_t)/(H_s)_{\star_t}$,
which $T_{-t}$ identifies with
$\mathcal E/(H_{s-t})$; its support
flag map is unchanged.
Every parity receiver and saturation
argument (HFD18)--(HFD22) remains exact.

For the same physical divisor at time
$s$ in the transported presentation,
the generator is $h_{s+t}=U_t h_s$.
In that notation the exact identification is
\begin{equation}\tag{HFD31}
 \mathcal D^{\mathrm{flag},\star_t}_{s+t,L}
       \simeq\mathcal D^{\mathrm{flag}}_{s,L}.
\end{equation}
Ordinary evaluation on the unchanged
amplitude $f$ is not silently substituted
for the transported amplitude map.
The ordinary and transported constructions,
their full coefficient quotients, their
off-real jet receivers and their mixed
prime therefore have specified connecting
morphisms throughout.
